\scanpage{20}\mspage{6 - 44 -}\hfill n° 392

est principal. Inversement, pour un nombre premier $p$, considérer le sous-anneau $C=Z+Bp$ de $B$ ; on a $\underline f=\underline f_{C/B}=Bp$, et $C/\underline f$ est isomorphe à $\underline Z/p\underline Z=\underline F_p$ ; observer enfin que $C/\underline f$ est un $C$-module isomorphe à $C^*/B^*$, donc à $B/C$, et en conclure que $B$ est nécessairement un $\underline Z$-module de dimension 2).

6) Les hypothèses sur $A$, $B$, $C$ étant celles de l'exerc. 2, soient $\underline p$ un idéal premier de $A$ de hauteur 1, $\underline P$ un idéal premier de $C$ au-dessus de $\underline p$. Montrer que les conditions suivantes sont équivalentes :

a) $\underline P_{\underline p}$ ne contient pas $(\underline D_{C/A})_{\underline p}$.

b) $\underline P_{\underline p}$ est une des composantes primaires de $\underline p C_{\underline p}$ dans $C_{\underline p}$, et le corps des fractions de $C/\underline P$ est une extension séparable de celui de $A/\underline p$.

c) Il existe un seul idéal premier $\underline P'$ de $B$ au-dessus de $\underline P$, $\underline P'$ est non ramifié sur $\underline p$, et les corps des fractions de $C/\underline P$ et de $B/\underline P'$ sont identiques.

(Se ramener au cas où $A$ est un anneau de valuation discrète. Pour voir que a) entraîne b), utiliser l'exerc. 3 b), l'exerc. 1 b) pour évaluer la longueur de $C/\underline Q$, où $\underline Q$ est la composante $\underline P$-primaire de $C\underline p$, et enfin le chap. V, \S{} 1, exerc. 16. Utiliser encore l'exerc. 1 b) pour voir que b) entraîne c). Enfin, pour montrer que c) implique a), considérer les complétés des anneaux semi-locaux $B$ et $C$).

7)\quad a) Sous les hypothèses de la prop. \quad, montrer que si les $x^i$, $(0\leq i\leq n-1)$, forment une base du $A$-module $B$, il en est de même des $y^i$, $(0\leq i\leq n-1)$, pour tout $y\in B$ suffisamment voisin de $x$.

b) Soient $A$ un anneau de valuation discrète, $K$ son corps des fractions, $B$ la fermeture algébrique de $A$ dans une extension finie séparable $L$ de $K$. Soit $\underline P$ un idéal premier de $B$ tel que le corps des fractions de $B/\underline P$ soit une extension séparable du corps résiduel $\bar K$ de $A$. Montrer qu'il existe $x\in B$ tel que $L=K(x)$ et tel que le conducteur $f_{B/A[x]}$ de $B$ dans $A[x]$ soit étranger à $\underline P$.

c) Les hypothèses sur $A$ et $B$ étant celles de b), supposons qu'il existe $n=[L:K]$ idéaux premiers distincts de $B$ au-dessus de l'idéal maximal $\underline p$ de $A$, (on dit alors que $B$ est ``complètement décomposé''). Montrer que, pour qu'il existe $x\in B$ tel que $B=A[x]$, il faut et il suffit que l'on ait $n\leq\operatorname{Card}(\bar K)$, (utiliser le lemme de Hensel et le th. d'approximation du chap. VI, \S{} 7, n° \quad).
% archive sheet number: 19
